% -*- coding: utf-8 -*-
%-------------------------designed by zcf--------------
\documentclass[UTF8,a4paper,10pt,fontset=fandol]{ctexart}
\usepackage[left=3.17cm, right=3.17cm, top=2.74cm, bottom=2.74cm]{geometry}
\usepackage{amsmath}
\usepackage{graphicx,subfig}
\usepackage{float}
\usepackage{cite}
\usepackage{caption}
\usepackage{enumerate}
\usepackage{booktabs} %表格
\usepackage{multirow}
\newcommand{\tabincell}[2]{\begin{tabular}{@{}#1@{}}#2\end{tabular}}  %表格强制换行
%-------------------------字体设置--------------
% \usepackage{times} 
\usepackage{ctex}
% 使用 TeX Live 自带的 Fandol 中文字体，避免依赖系统字体。
\newcommand{\yihao}{\fontsize{26pt}{36pt}\selectfont}           % 一号, 1.4 倍行距
\newcommand{\erhao}{\fontsize{22pt}{28pt}\selectfont}          % 二号, 1.25倍行距
\newcommand{\xiaoer}{\fontsize{18pt}{18pt}\selectfont}          % 小二, 单倍行距
\newcommand{\sanhao}{\fontsize{16pt}{24pt}\selectfont}  %三号字
\newcommand{\xiaosan}{\fontsize{15pt}{22pt}\selectfont}        % 小三, 1.5倍行距
\newcommand{\sihao}{\fontsize{14pt}{21pt}\selectfont}            % 四号, 1.5 倍行距
\newcommand{\banxiaosi}{\fontsize{13pt}{19.5pt}\selectfont}    % 半小四, 1.5倍行距
\newcommand{\xiaosi}{\fontsize{12pt}{18pt}\selectfont}            % 小四, 1.5倍行距
\newcommand{\dawuhao}{\fontsize{11pt}{11pt}\selectfont}       % 大五号, 单倍行距
\newcommand{\wuhao}{\fontsize{10.5pt}{15.75pt}\selectfont}    % 五号, 单倍行距
%-------------------------章节名----------------
\usepackage{ctexcap} 
\CTEXsetup[name={,、},number={ \chinese{section}}]{section}
\CTEXsetup[name={（,）},number={\chinese{subsection}}]{subsection}
\CTEXsetup[name={,.},number={\arabic{subsubsection}}]{subsubsection}
%-------------------------页眉页脚--------------
\usepackage{fancyhdr}
\pagestyle{fancy}
\lhead{\kaishu \leftmark}
% \chead{}
\rhead{\kaishu 操作系统}%加粗\bfseries
\lfoot{}
\cfoot{\thepage}
\rfoot{}
\renewcommand{\headrulewidth}{0.1pt}  
\renewcommand{\footrulewidth}{0pt}%去掉横线
\newcommand{\HRule}{\rule{\linewidth}{0.5mm}}%标题横线
\newcommand{\HRulegrossa}{\rule{\linewidth}{1.2mm}}
%-----------------------伪代码------------------
\usepackage{algorithm}  
\usepackage{algorithmicx}  
\usepackage{algpseudocode}  
\floatname{algorithm}{Algorithm}  
\renewcommand{\algorithmicrequire}{\textbf{Input:}}  
\renewcommand{\algorithmicensure}{\textbf{Output:}} 
\usepackage{lipsum}  
\makeatletter
\newenvironment{breakablealgorithm}
  {% \begin{breakablealgorithm}
  \begin{center}
     \refstepcounter{algorithm}% New algorithm
     \hrule height.8pt depth0pt \kern2pt% \@fs@pre for \@fs@ruled
     \renewcommand{\caption}[2][\relax]{% Make a new \caption
      {\raggedright\textbf{\ALG@name~\thealgorithm} ##2\par}%
      \ifx\relax##1\relax % #1 is \relax
         \addcontentsline{loa}{algorithm}{\protect\numberline{\thealgorithm}##2}%
      \else % #1 is not \relax
         \addcontentsline{loa}{algorithm}{\protect\numberline{\thealgorithm}##1}%
      \fi
      \kern2pt\hrule\kern2pt
     }
  }{% \end{breakablealgorithm}
     \kern2pt\hrule\relax% \@fs@post for \@fs@ruled
  \end{center}
  }
\makeatother
%------------------------代码-------------------
\usepackage{xcolor} 
\usepackage{listings} 
\lstset{
breaklines,%自动换行
basicstyle=\small,
escapeinside=``,
keywordstyle=\color{ blue!70} \bfseries,
commentstyle=\color{red!50!green!50!blue!50},% 
stringstyle=\ttfamily,% 
extendedchars=false,% 
linewidth=\textwidth,% 
numbers=left,% 
numberstyle=\tiny \color{blue!50},% 
frame=trbl% 
rulesepcolor= \color{ red!20!green!20!blue!20} 
}
%------------超链接----------
\usepackage[colorlinks,linkcolor=black,anchorcolor=blue]{hyperref}
%------------------------TODO-------------------
\usepackage{enumitem,amssymb}
\newlist{todolist}{itemize}{2}
\setlist[todolist]{label=$\square$}
% for check symbol 
\usepackage{pifont}
\newcommand{\cmark}{\ding{51}}%
\newcommand{\xmark}{\ding{55}}%
\newcommand{\done}{\rlap{$\square$}{\raisebox{2pt}{\large\hspace{1pt}\cmark}}\hspace{-2.5pt}}
\newcommand{\wontfix}{\rlap{$\square$}{\large\hspace{1pt}\xmark}}
%------------------------水印-------------------
\usepackage{tikz}
\usepackage{xcolor}
\usepackage{eso-pic}

\newcommand{\watermark}[3]{\AddToShipoutPictureBG{
\parbox[b][\paperheight]{\paperwidth}{
\vfill%
\centering%
\tikz[remember picture, overlay]%
  \node [rotate = #1, scale = #2] at (current page.center)%
    {\textcolor{gray!80!cyan!30!magenta!30}{#3}};
\vfill}}}



%———————————————————————————————————————————正文———————————————————————————————————————————————
%----------------------------------------------
\begin{document}
\begin{titlepage}
    \begin{center}
    \includegraphics[width=0.8\textwidth]{NKU.png}\\[1cm]    
    \textsc{\Huge \kaishu{\textbf{南\ \ \ \ \ \ 开\ \ \ \ \ \ 大\ \ \ \ \ \ 学}} }\\[0.9cm]
    \textsc{\huge \kaishu{\textbf{计\ \ 算\ \ 机\ \ 学\ \ 院}}}\\[0.5cm]
    \textsc{\Large \textbf{操作系统}}\\[0.8cm]
    \HRule \\[0.9cm]
    \vspace{0.4cm}
    \HRule \\[2.0cm]
    \centering
    \textsc{\LARGE \kaishu{白景仁\ \ 陈彦兴\ \ 何闻硕}}\\[0.5cm]
    \textsc{\LARGE \kaishu{年级\ :\ 2026级}}\\[0.5cm]
    \textsc{\LARGE \kaishu{专业\ :\ 计算机科学与技术}}\\[0.5cm]
    \textsc{\LARGE \kaishu{指导教师\ : 袁晓洁\ \ 李浩然}}\\[0.5cm]
    \vfill
    {\Large \today}
    \end{center}
\end{titlepage}
%-------------摘------要--------------
\newpage
\thispagestyle{empty}
\renewcommand{\abstractname}{\kaishu \sihao \textbf{摘要}}
    \begin{abstract}
        \noindent  %顶格
        \textbf{\\\ 关键字：}\textbf{} \\\ \\\
    \end{abstract}
%----------------------------------------------------------------
\tableofcontents
%----------------------------------------------------------------
\newpage
\watermark{60}{10}{NKU}
\setcounter{page}{1}
\section{实验基本信息与小组分工}
\subsection{实验名称与完成日期}
Lab1：最小可执行内核。完成日期：待填写。
\subsection{练习分工}
待填写各成员负责的练习与调试任务。
\subsection{报告分工}
待填写各成员负责的报告内容。

\section{实验目的}
待填写本实验的学习目标。

\section{实验环境}
\subsection{硬件、操作系统与工具链}
待填写实际环境及 QEMU、GDB、RISC-V 工具链版本。
\subsection{AI 工具与模型}
待填写本实验实际使用的工具、模型；未记录的版本如实注明。

\section{实验整体逻辑分析}
\subsection{本章节的逻辑主线}
待说明复位代码、OpenSBI、内核入口、初始化与输出之间的关系。
\subsection{功能的逐步执行与核心模块}
待结合链接脚本、入口汇编、初始化函数和输出模块说明执行顺序。

\section{实验内容与练习解答}
\subsection{练习一：理解内核启动中的程序入口操作}

\subsubsection{\texttt{la sp, bootstacktop} 的操作与目的}

\texttt{la} 是一条加载地址的伪指令。\texttt{la sp, bootstacktop} 会将符号
\texttt{bootstacktop} 的地址加载到栈指针寄存器 \texttt{sp} 中。
\texttt{bootstacktop} 表示启动栈的栈顶地址。设置 \texttt{sp} 的目的是为后续的
函数调用、返回地址保存以及局部变量分配建立栈空间。

在本实验的 \texttt{entry.S} 中，启动栈通过下面的定义预留：

\begin{verbatim}
bootstack:
    .space KSTACKSIZE
bootstacktop:
\end{verbatim}

当前 \texttt{KSTACKSIZE} $= 2 \times 4096 = 8192$ 字节，因此
\texttt{bootstacktop = bootstack + KSTACKSIZE}。启动栈位于内核的
\texttt{.data} 段中，其最终地址由 \texttt{tools/kernel.ld} 链接脚本决定，
而不是简单固定为 \texttt{0x80000000}。链接脚本将内核的基地址设置为
\texttt{0x80200000}，并对数据段进行页对齐。

需要区分 QEMU 启动过程中的几个地址：CPU 复位后从
\texttt{0x1000} 开始执行，随后进入位于约 \texttt{0x80000000} 的 OpenSBI；
QEMU 将内核映像加载到 \texttt{0x80200000}；OpenSBI 完成初始化后，把控制权交给
内核入口 \texttt{kern\_entry}。

\subsubsection{\texttt{tail kern\_init} 的操作与目的}

\texttt{tail kern\_init} 将控制权转移到 C 语言函数 \texttt{kern\_init}，
开始执行内核初始化代码。该指令采用尾调用形式，不再为返回到
\texttt{entry.S} 额外保存新的返回地址。由于内核初始化完成后会进入死循环，
程序不需要返回到 \texttt{entry.S}，因此使用尾调用可以简化内核启动过程。

\subsection{练习二：使用 GDB 验证启动流程}
以下摘录本次调试中已获得的有效输出，按启动阶段组织。

\subsubsection{复位地址与最初六条指令}
\begin{verbatim}
(gdb) x/10i 0x1000
   0x1000: auipc t0,0x0
=> 0x1004: addi  a2,t0,40
   0x1008: csrr  a0,mhartid
   0x100c: ld    a1,32(t0)
   0x1010: ld    t0,24(t0)
   0x1014: jr    t0
\end{verbatim}
这里只保留六条启动指令；原输出从 \texttt{0x1018} 开始是将数据
按指令解码的结果。箭头表明采集时已执行第一条指令，暂停在 \texttt{0x1004}。
\begin{verbatim}
(gdb) x/2gx 0x1018
0x1018: 0x0000000080000000  0x0000000087e00000
\end{verbatim}
第一条指令使 \texttt{t0 = 0x1000}；随后设置
\texttt{a2 = 0x1028}，读取 hart ID 到 \texttt{a0}，并将设备树地址
\texttt{0x87e00000} 读入 \texttt{a1}。
\texttt{ld t0,24(t0)} 的读取地址为
\texttt{0x1000 + 24 = 0x1018}，读到 \texttt{0x80000000} 后，
\texttt{jr t0} 跳转到 OpenSBI。

\subsubsection{OpenSBI 入口寄存器与指令}
在 \texttt{0x80000000} 断点处，寄存器输出的关键部分如下：
\begin{verbatim}
ra  0x0         0x0
sp  0x0         0x0
t0  0x80000000  2147483648
a0  0x0         0
a1  0x87e00000  2279604224
a2  0x1028      4136
pc  0x80000000  0x80000000
\end{verbatim}
入口反汇编摘录如下：
\begin{verbatim}
(gdb) x/30i 0x80000000
   0x80000000: add s0,a0,zero
   0x80000004: add s1,a1,zero
   0x80000008: add s2,a2,zero
   0x8000000c: jal 0x80000678
   0x80000010: add a6,a0,zero
   0x80000014: add a0,s0,zero
   0x80000018: add a1,s1,zero
   0x8000001c: add a2,s2,zero
\end{verbatim}
这些指令属于 OpenSBI，尚未进入本实验的内核入口。

\subsubsection{内核入口与初始化指令}
\begin{verbatim}
(gdb) x/10i $pc
=> 0x80200000 <kern_entry>:    auipc sp,0x3
   0x80200004 <kern_entry+4>:  mv    sp,sp
   0x80200008 <kern_entry+8>:  j     0x8020000a <kern_init>
   0x8020000a <kern_init>:     auipc a0,0x3
   0x8020000e <kern_init+4>:   addi  a0,a0,-2
   0x80200012 <kern_init+8>:   auipc a2,0x3
   0x80200016 <kern_init+12>:  addi  a2,a2,-10
   0x8020001a <kern_init+16>:  addi  sp,sp,-16
   0x8020001c <kern_init+18>:  li    a1,0
   0x8020001e <kern_init+20>:  sub   a2,a2,a0
\end{verbatim}
PC 已到达 \texttt{0x80200000}，验证控制权交给了内核。
入口前两条指令将栈指针设置为 \texttt{0x80203000}，随后跳到
\texttt{kern\_init}。初始化指令准备 BSS 清零的参数；从此反汇编计算得到
\texttt{a0}、\texttt{a2} 最初都为 \texttt{0x80203008}，相减后长度为零。
这些数值由指令推算，未另行采集执行后的寄存器输出。

\subsection{补充观察：输出调用与死循环}
\subsubsection{cprintf 调用与返回地址}
\begin{verbatim}
0x0000000080200036 <+44>: 01e000ef jal 0x80200054 <cprintf>
0x000000008020003a <+48>: a001     j   0x8020003a <kern_init+48>
\end{verbatim}
\texttt{jal} 调用 \texttt{cprintf} 并将下一条指令地址
\texttt{0x8020003a} 写入 \texttt{ra}。文本输出在被调用函数及其底层输出函数中完成。

\subsubsection{进入固件后无法使用 finish}
\begin{verbatim}
(gdb) finish
"finish" not meaningful in the outermost frame.
(gdb) bt
#0  0x000000008000caae in ?? ()
\end{verbatim}
此时地址位于 OpenSBI 范围，GDB 未能识别完整调用栈，不能据此认定
当前函数确实是程序最外层函数。

\subsubsection{输出完成后的死循环}
输出完成后，反汇编的有效摘录如下：
\begin{verbatim}
(gdb) x/10i $pc
=> 0x8020003a <kern_init+48>: j    0x8020003a <kern_init+48>
   0x8020003c <cputch>:      addi sp,sp,-32
   0x8020003e <cputch+2>:    sd   ra,24(sp)
   0x80200040 <cputch+4>:    sd   a1,8(sp)
   0x80200042 <cputch+6>:    jal  0x80200088 <cons_putc>
\end{verbatim}
当前指令跳回自身，对应 \texttt{while (1)}。
紧邻的 \texttt{cputch} 是独立函数，内存地址相邻不意味着死循环后会执行它。

\subsection{AI 辅助学习与调试过程}
本实验未进行代码生成或功能修改，功能实现提示词与代码实现迭代不适用。
待如实整理 AI 的实际用途、遇到的误判以及通过源码和调试结果核验的过程；
不补造原始提示词或对话记录。

\section{测试与验证}
\subsection{编译与运行结果}
待放入实际 \texttt{make}、\texttt{make qemu} 输出或截图。
\subsection{启动流程的 GDB 验证结果}
待汇总复位地址、OpenSBI 入口、内核入口、栈指针变化和最终输出证据。
\subsection{验证范围与未执行项目}
待说明实际验证范围；未运行评分脚本时，不填写 \texttt{make grade} 通过。

\section{实验总结与收获}
\subsection{实验知识点与操作系统原理的联系和差异}
待整理启动流程、内存布局、栈、特权级与固件接口等知识点。
\subsection{重要但本实验尚未涉及的操作系统知识点}
待说明进程调度、虚拟内存、文件系统等尚未涉及的内容。
\subsection{调试中遇到的问题与解决方法}
待整理地址混淆、数据误反汇编、OpenSBI 中无法回溯等实际问题。
\subsection{AI 协作经验与个人收获}
待填写实际经验与学习收获。

%----------------------------------------------------------------
\newpage
\bibliographystyle{plain}
\bibliography{references} 
\end{document}
